% Part: turing-machines
% Chapter: undecidability
% Section: unsolvability-decision-problem

\documentclass[../../../include/open-logic-section]{subfiles}

\begin{document}

\olfileid{tur}{und}{dec}
\olsection{निर्णय समस्या}

दिलेल्या !!{sentence} वाक्याची वैधता ठरवण्याची
प्रभावी पद्धत असेल तेव्हा आणि केवळ तेव्हाच
प्रथम-क्रम तर्कशास्त्र \emph{निर्णेय} आहे,
असे आपण म्हणतो. प्रत्यक्षात अशी पद्धत नाही:
प्रथम-क्रम वाक्यांच्या वैधतेचा निर्णय
करण्याची समस्या सोडवता येत नाही.

हा महत्त्वाचा निषेधात्मक निष्कर्ष स्थापित
करण्यासाठी, निर्णय समस्या ट्यूरिंग यंत्राने
सोडवता येत नाही हे आपण सिद्ध करतो.
म्हणजे, फितीवर प्रथम-क्रम !!{sentence}
वाक्य असताना सुरू केले की ते !!{sentence} वाक्य वैध
आहे की नाही यानुसार अखेरीस थांबून
$1$ किंवा~$0$ फलित देईल, असे कोणतेही
ट्यूरिंग यंत्र नाही हे दाखवतो.
चर्च--ट्यूरिंग प्रबंधानुसार, संगणनक्षम
असलेले प्रत्येक फलन ट्यूरिंग-संगणनक्षम
असते. म्हणून हे “वैधतेचे फलन”
मुळात प्रभावी रीतीने संगणित करता
येत असेल, तर ते ट्यूरिंग-संगणनक्षमही
असेल. ते ट्यूरिंग-संगणनक्षम नसेल,
तर प्रभावी रीतीनेही संगणित करता
येणार नाही.

निर्णय समस्या सोडवता येत नाही हे
सिद्ध करण्यासाठी आपण थांबण्याच्या
समस्येचे तिच्याकडे न्यूनीकरण करू.
याचा अर्थ पुढीलप्रमाणे. ट्यूरिंग
यंत्राचे वर्णन~$e$ यांनी होत असेल,
तर ते यंत्र आदान~$w$ वर थांबते
तेव्हा~$1$ आणि अन्यथा~$0$ देणारे
फलन~$h(e,w)$ ट्यूरिंग-संगणनक्षम
नाही, हे आपण सिद्ध केले आहे.
प्रथम-क्रम वाक्यांची वैधता ठरवणारे
ट्यूरिंग यंत्र असेल, तर~$h$ संगणित
करणारे ट्यूरिंग यंत्रही असेल, हे
आपण दाखवू. $h$ ट्यूरिंग यंत्राने
संगणित करता येत नसल्यामुळे वैधतेचा
निर्णय करणारे ट्यूरिंग यंत्रही असू
शकत नाही.

या युक्तीतील पहिली पायरी अशी आहे:
प्रत्येक आदान~$w$ आणि ट्यूरिंग
यंत्र~$M$ साठी, $M$ चा सूचनासंच
आणि आदान~$w$ दर्शवणारे
!!a{sentence} वाक्य $!T(M, w)$,
तसेच “$M$ अखेरीस थांबते”
हे व्यक्त करणारे !!a{sentence}
वाक्य~$!E(M, w)$ प्रभावी रीतीने
लिहिता येतात, असे दाखवायचे;
आणि ती अशी असतील की
\begin{quote}
  $\Entails !T(M, w) \lif !E(M,w)$ हे तेव्हाच खरे, जेव्हा $M$ हे आदान~$w$ वर थांबते.
\end{quote}
आपल्या पुराव्याचा मुख्य भाग ही
वाक्ये $!T(M, w)$ आणि~$!E(M, w)$
यांचे वर्णन करण्यात आणि
$!T(M, w) \lif !E(M, w)$ हे
आदान~$w$ वर~$M$ थांबते तेव्हाच
वैध असते हे पडताळण्यात जाईल.

\end{document}
